\section[Electronic Lab Notebook]{Electronic Lab Notebook (ELN)}\label{ssc:ELN}

\subsection{Allgemeine Informationen}

Electronic Lab Notebooks ermöglichen es Forschende, experimentelle
Verfahren, Protokolle, Notizen und Daten über ihren Computer oder ein mobiles
Gerät zu organisieren und zu speichern. ELNs bieten gegenüber dem traditionellen
Papiernotizbuch mehrere Vorteile bei der Dokumentation von Forschungsarbeiten in
der aktiven Phase eines Projekts. Dazu gehören die Durchsuchbarkeit innerhalb
und zwischen den Notizbüchern, die sichere Speicherung mit mehreren Redundanzen,
der Fernzugriff auf die Notizbücher und die Möglichkeit, die Beschreibungen von
Experimenten oder Simulationen problemlos mit anderen Teammitgliedern und
Mitarbeitenden zu teilen.

eLabFTW \includegraphics[height=9pt]{eLabFTW-logo3-ohneText.png} ist ein OpenSource, generisches, browserbasiertes ELN, entwickelt/initialisiert von Nicolas Carpi (Institut Pasteuer, Paris) im Jahr 2012. Die Daten werden in MySQL/MariaDB gespeichert. Es wird von vielen Entwickler*innen gepflegt und ist weltweit im Einsatz (Berkley, Indian Institute of Technology, KIT, ...).
\begin{itemize}
  \item Zugang: \url{https://elabftw.iwt.zz/}
  \item Frei definierbarer Status der Experimente ("beendet", "läuft", ...)
  \item Definition von Vorlagen für eine schrittweise Ablaufbeschreibung von
        Messungen
  \item Jedes Experiment erhält eine eindeutige ID und kann leicht als pdf-Datei
        zusammengefasst werden
  \item Frei definierbare Kategorien/Tags
  \item Grafischer Texteditor zur Beschreibung Ihrer Experimente und Simulationen
  \item Dateianhänge von Daten mit Vorschau der gängigen Formate (pdf, tiff,
        png, ...)
  \item Verknüpfung von Experimenten
  \item Zeitstempel-Dienst (RFC 3161, z.B. DFN)
  \item Datenimport/-export (csv, zip, json, ...)
  \item Zugriff mit z.B. Python über eine Anwendungsprogrammierschnittstelle (API)
\end{itemize}



\subsection{IWT-Vorgaben zur Dokumentation und Struktur}

Um die wissenschaftlichen Experimente übersichtlich, klar und verständlich für andere zu dokumentieren und auch noch Jahre später schnell auffinden zu können, werden die Titel in eLabFTW einheitlich mit dieser Systematik aus sechs Komponenten vorgegeben: \\
 \\
\{Versuchsnummer\}, \{Probename/Kürzel\}, Material, \{Gerät/Anlage\}, Methode, Projektkürzel \\
 \\
Die Angaben in eckigen Klammern sind optional, aber die drei Angaben für ``Material'', ``Methode'' und das ``Projekt'' bzw. ``Projektkürzel'' sind zwingend erforderlich, es können aber auch mehr als die oben genannten sein. Mit genau derselben Systematik werden Tags (d.h. Marker bzw. Schlagwörter) vergeben. Der Grund für diese Redundanz liegt in dem SQL basierten Datenbank-System von eLabFTW, das mithilfe von Tags wesentlich schnellere Antworten auf Suchanfragen liefert.

Zur Dokumentation der Ergebnisse kann für ein Experiment (was auch eine Simulation sein kann) eine PDF-Datei erstellt werden, indem Sie im Bearbeitungsmodus einfach auf 'PDF erstellen' klicken (siehe
\texttt{‘04\_PrimaryData’} in Tabelle \ref{table:paper-directory-structure}).
Diese Art der Dokumentation enthält einen eindeutigen QR-Code und ist für Daten, die veröffentlicht werden, erforderlich! Der QR-Code ermöglicht einen direkten Link zu den Daten zusammen mit einer eindeutigen ID.

Die \textbf{maximale Größe} der hochgeladenen Daten sollte \textbf{100Mbyte pro Datei} nicht überschreiten. Weitere Tipps und Tricks finden Sie im eLab-Dokument (bitte fügen Sie dort Kommentare hinzu, wenn Sie neue Möglichkeiten der Nutzung entdecken) auf dem IWT file server   \\ ($\rightarrow$Austausch $\rightarrow$Forschungsdatenmanagement $\rightarrow$eLabFTW)

Eine Vorlage für den in jedem Experiment vorhandenen JSON-Feld mit Angaben zu den Metadaten (siehe Kapitel \ref{ssc:Doku+Metadata}) ist ebenfalls auf dem IWT-Fileserver abgelegt.
